\documentclass[11pt]{article}

\usepackage[utf8]{inputenc}
\usepackage[T1]{fontenc}
\usepackage{lmodern}
\usepackage{geometry}
\usepackage{amsmath,amssymb,amsfonts,amsthm,mathtools}
\usepackage{bm}
\usepackage{enumitem}
\usepackage{microtype}
\usepackage{hyperref}
\usepackage{titlesec}
\usepackage{booktabs}
\usepackage{array}
\usepackage{setspace}
\usepackage{mathrsfs}
\usepackage{physics}

\geometry{margin=1in}
\hypersetup{colorlinks=true, linkcolor=black, urlcolor=blue, citecolor=black}
\setlist[enumerate]{topsep=0.4em,itemsep=0.35em}
\setlist[itemize]{topsep=0.3em,itemsep=0.25em}
\titleformat{\section}{\large\bfseries}{\thesection.}{0.5em}{}
\titleformat{\subsection}{\normalsize\bfseries}{\thesubsection.}{0.5em}{}
\setstretch{1.06}

\newcommand{\Aether}{\AE{}ther}
\newcommand{\Aflow}{\AE{}ther-flow}
\newcommand{\TheoryName}{\Aflow{} Interpretation of Relativity}
\newcommand{\TheoryProgram}{\Aether{} / \Aflow{} framework}
\newcommand{\CompletedMetric}{g^{\sharp}}

\newcommand{\ProjectionOperatorCore}{\mathfrak S^{\mathrm{disp,resp,projop}}}
\newcommand{\CurrentBodyImageCore}{\mathfrak S^{\mathrm{disp,resp,cbimg}}}
\newcommand{\CanonicalImageBody}[1]{\mathcal U_{#1}^{\mathrm{disp,resp,cbimg}}}

\newtheorem{definition}{Definition}
\newtheorem{proposition}{Proposition}
\newtheorem{remark}{Remark}
\newtheorem{corollary}{Corollary}
\newtheorem{theorem}{Theorem}

\title{\textbf{The \TheoryName{}}\\[0.4em]
\large Primitive-Reservoir Observer-Reduction\\
\large Section-Centered Hidden-Response\\
\large Canonical Image-Body Necessity / Uniqueness\\
\large Next Benchmark Criterion}
\author{}
\date{}

\begin{document}

\maketitle

\begin{abstract}
The current active-primary theorem now proves that the canonical image body \(\CanonicalImageBody{\sigma}\) is already the unique benchmark-facing body compatible with the fixed current displacement body, the fixed surjective projection, the fixed coercive response operator, the fixed exact-shell image, and the fixed current body-response route. Another same-output deeper-origin relay that merely preserves that same body and therefore merely reproduces the same canonical current-body image core \(\CurrentBodyImageCore\) no longer counts as the honest default next benchmark-facing manuscript.

This note records the routing consequence of that gain. The next honest benchmark-facing move must now derive or justify the remaining nontrivial constituents or relations inside the current image-core frontier from still more primitive explicit substrate structure in a way that changes benchmark-facing status, or derive the full current image core itself if that deeper package is genuinely cleaner. The benchmark boundary remains fixed throughout. The one operative metric is still exactly \(\CompletedMetric\), the ordinary matter route is unchanged, there is no second operative metric, the low-energy matter law is not enlarged, and no stronger promotion language is licensed. The positive result is therefore routing discipline below the canonical-image-body necessity / uniqueness frontier.
\end{abstract}

\section{Introduction}

The active derivational program of \TheoryName{} is no longer merely at current-body-image-core scope in the broad sense recorded by the prior collapse theorem. The preceding canonical-image-body necessity / uniqueness theorem already showed that, on the recorded section-centered hidden-response route, the body constituent \(\CanonicalImageBody{\sigma}\) carries no remaining benchmark-facing presentational freedom once the fixed current displacement body and the fixed projection / response data are held fixed. That theorem therefore spent the body-level forcing or body-level uniqueness option that had remained open inside the prior image-core frontier.

The present note records the routing consequence of that gain. It does not reopen the larger full substrate-body presentation, the larger projection-operator core \(\ProjectionOperatorCore\), the larger full displacement-response substrate law, or the already-spent body constituent inside \(\CurrentBodyImageCore\) as if any of those were again independently live. It records only which kind of next manuscript would now count as an honest benchmark-facing gain once the body constituent is no longer the open burden.

\paragraph{Framework and claim boundary.}
\TheoryProgram{} denotes the broader \Aether{} / \Aflow{} framework built on the claims that the \Aether{} is the underlying four-dimensional substrate of reality, the \Aflow{} is its intrinsic ordered motion, observed three-dimensional space is the local experiential slice of that deeper substrate, S-time is the experienced order of change arising from matter, light, and the \Aflow{}, and observed expansion is the three-dimensional appearance of deeper four-dimensional ordered motion. Within that framework, \TheoryName{} denotes the exact-closure theory in which the effective gravitational dynamics are adopted to be exactly Einsteinian with universal matter coupling. In the active sequence, \emph{adoption} means use of that established relativistic dynamics without claiming substrate derivation, while \emph{derivation} is reserved for a first-principles recovery from explicit substrate variables.

The present note stays strictly inside that boundary. It records only which kind of next manuscript would now count as an honest benchmark-facing gain below the canonical-image-body necessity / uniqueness frontier.

\section{Fixed Inputs}

\begin{definition}[Canonical-image-body necessity / uniqueness frontier inputs]
The criterion recorded here uses the following fixed inputs.
\begin{enumerate}
    \item The benchmark exact-closure package which fixes the public one-metric GR target of the repository.
    \item The benchmark gatekeeping layer including the recorded safeguard that blocks same-output deeper-origin relays from counting as new front-line progress once the live benchmark-relevant datum is unchanged.
    \item The earlier theorem showing that the full projection-operator-core body presentation collapses benchmark-facingly to the smaller canonical current-body image core \(\CurrentBodyImageCore\).
    \item The later theorem proving that the canonical image body \(\CanonicalImageBody{\sigma}\) is already the unique benchmark-facing body compatible with the fixed current route data inside that smaller image core.
\end{enumerate}
\end{definition}

\begin{remark}
The present note does not replace those manuscripts. It records the routing consequence of reading them together under the active safeguard against recursive depth drift.
\end{remark}

\section{Why the Live Burden Now Sits Below Body Restatement}

\begin{proposition}[The live burden is renewed below canonical-image-body restatement]
At the current stage of the primitive-reservoir observer-reduction line, the canonical-image-body necessity / uniqueness theorem renews the live benchmark-facing burden below any mere restatement of the canonical image body \(\CanonicalImageBody{\sigma}\).
\end{proposition}

\begin{proof}
That theorem changes benchmark-facing status because it proves that the present image body is already uniquely forced on the recorded route once the fixed current displacement body and the fixed projection / response data are held fixed. Once that uniqueness statement is in hand another manuscript that merely preserves the same canonical image body no longer removes any remaining benchmark-facing freedom. The live burden therefore no longer sits on the task of restating or relocating \(\CanonicalImageBody{\sigma}\) itself. It sits on deriving or sharpening the remaining nontrivial constituents or relations inside the current image-core data.
\end{proof}

\begin{definition}[Benchmark-neutral same-body relay below the canonical-image-body necessity / uniqueness frontier]
A \emph{benchmark-neutral same-body relay below the canonical-image-body necessity / uniqueness frontier} is a manuscript that preserves the same benchmark one-metric output, the same ordinary matter route, the same canonical image body \(\CanonicalImageBody{\sigma}\), the same canonical current-body image core \(\CurrentBodyImageCore\), and the same downstream benchmark-facing consequences while merely relocating the unexplained substrate layer deeper, renaming the same body constituent, restoring the already-screened larger full substrate-body presentation, or re-expanding back to the larger projection-operator core \(\ProjectionOperatorCore\), the full displacement-response substrate law, or the larger pullback-body-operator, pullback-operator, pullback-quadratic, hidden-response, residual-normal-form, or pre-proto-section presentations without a new derivation, new forcing theorem, or comparably sharp new gain.
\end{definition}

\section{The Current Post-Uniqueness Criterion}

\begin{definition}[Admissible next benchmark-facing manuscript below the canonical-image-body necessity / uniqueness frontier]
Below the current canonical-image-body necessity / uniqueness frontier a manuscript counts as an admissible next benchmark-facing move only if it does at least one of the following:
\begin{enumerate}
    \item derives or justifies the restricted projection / canonical-lift relation on the fixed current displacement body from still more primitive explicit substrate structure in a genuinely benchmark-facing way;
    \item derives or justifies the exact-shell image or the response operator restricted to the canonical image from still more primitive explicit substrate structure in a way that sharpens the present route;
    \item derives or justifies the full canonical current-body image core \(\CurrentBodyImageCore\) in one cleaner deeper package if that unified derivation is genuinely available;
    \item proves a sharper necessity, uniqueness, or compatibility theorem for some remaining nontrivial constituent or relation inside \(\CurrentBodyImageCore\);
    \item does one of the previous four while preserving the same benchmark one-metric package, the same ordinary matter route, and the same claim boundary.
\end{enumerate}
\end{definition}

\begin{theorem}[Next honest benchmark-facing task below the canonical-image-body necessity / uniqueness frontier]
The next honest benchmark-facing manuscript below the current canonical-image-body necessity / uniqueness frontier must derive or justify some remaining nontrivial constituent or relation inside the current image-core data from still more primitive explicit substrate structure in a benchmark-relevant way, or else derive the full canonical current-body image core \(\CurrentBodyImageCore\) in a genuinely cleaner deeper package. A manuscript that only repeats the same canonical image body through another deeper relay, restores the larger full substrate-body presentation, re-expands the discussion back to the larger projection-operator core \(\ProjectionOperatorCore\) or the full displacement-response substrate law, or replays the larger already-spent presentations without such a new gain does not satisfy the criterion.
\end{theorem}

\begin{proof}
By Proposition 1 the live burden already lies below any mere restatement of the canonical image body \(\CanonicalImageBody{\sigma}\). Therefore a subsequent manuscript must improve on the already-forced body constituent rather than merely accompany it. A deeper-origin manuscript can do so only if it changes benchmark-facing status by more than depth alone. If it simply reproduces the same body and its same downstream outputs the routing safeguard classifies it as benchmark neutral. The remaining honest alternatives are therefore exactly those stated in the definition: derive or justify a remaining constituent or relation inside the current image-core data or derive the full image core in a cleaner deeper package.
\end{proof}

\section{What the Next Move Should Not Be}

\begin{proposition}[Screened next moves below the canonical-image-body necessity / uniqueness frontier]
Below the current canonical-image-body necessity / uniqueness frontier none of the following is the default next benchmark-facing move:
\begin{enumerate}
    \item another same-output deeper-origin relay that merely preserves the current canonical image body \(\CanonicalImageBody{\sigma}\) and therefore merely reproduces the same current image core \(\CurrentBodyImageCore\) together with the same downstream benchmark-facing consequences;
    \item another restoration of the full substrate-body presentation as if the already-screened off-image directions were again independently live;
    \item another replay of the full projection-operator core \(\ProjectionOperatorCore\), the full displacement-response substrate law, or the larger pullback-body-operator, pullback-operator, pullback-quadratic, hidden-response, residual-normal-form, or pre-proto-section presentations as if those already-spent larger presentations were again the honest current burden;
    \item a manuscript that introduces a second operative metric, enlarges the ordinary matter law, or uses stronger promotion language.
\end{enumerate}
\end{proposition}

\begin{proof}
Item (1) fails the preceding criterion because it leaves the renewed live burden unchanged. Item (2) likewise fails because it restores a larger presentation whose off-image freedom is already screened out. Item (3) fails because it re-expands to larger already-spent presentations rather than deriving or sharpening the remaining live constituents or relations inside the current image-core data. Item (4) violates the fixed claim boundary of the benchmark package and therefore cannot count as a disciplined next step on the active line.
\end{proof}

\begin{remark}
This screening statement is not a denial that those deeper or alternative manuscripts may matter strategically. It is a routing statement: they are not the default way to spend the next benchmark-facing move below the canonical-image-body necessity / uniqueness frontier unless they do the same benchmark-facing work named above.
\end{remark}

\section{Conclusion}

The positive result of this note is routing discipline below the canonical-image-body necessity / uniqueness frontier. The present necessity / uniqueness theorem remains the live benchmark-facing gain because it already proves that the canonical image body is uniquely forced on the recorded route once the fixed current displacement body and the fixed projection / response data are held fixed. The earlier image-core collapse theorem remains preserved main-line history but under the recursive-depth safeguard it is no longer itself the default current frontier.

The scope remains narrow and honest. The benchmark package is unchanged: the one operative metric is still exactly \(\CompletedMetric\), the ordinary matter route is unchanged, there is no second operative metric, and the low-energy matter law is not enlarged. Nothing here upgrades adoption to derivation or licenses stronger promotion language.

The next open burden is therefore clear. The next benchmark-facing manuscript should derive or justify the restricted projection / canonical-lift relation, the exact-shell image, the response operator restricted to the canonical image, or the full current image core itself from still more primitive explicit substrate structure in a way that changes benchmark-facing status. Another same-output deeper relay that merely preserves the same body datum, another restoration of the full substrate-body presentation, another replay of the full projection-operator core or the full displacement-response substrate law, or another re-expansion back to the larger already-spent presentations is not the clean next manuscript target at current scope.

\end{document}
